\chapter{Responsible use of resources and an end-of-life strategy } 
The responsible use of resources and a well-defined end-of-life strategy are crucial to the project. Without a sustainable plan for decommissioning, the benefits of the project, such as advancing research in renewable energy and promoting food sustainability through small-scale greenhouses, could be overshadowed by the environmental impact of improperly discarded materials. 
During the design and construction phase of the project, opportunities to reuse previously used materials or use recycled material will be taken into consideration.
Once completed, the greenhouse can serve as a model to demonstrate the potential of engineering projects, inspire the adoption of renewable energy solutions, and promote education on food sustainability and self-sufficiency.
Once this post completion life cycle is over. The Project can be broken down and stripped of its parts. Some of the parts will go into storage to be reused for future engineering Thesis topics or other experiments to be used by the university at their own discretion. Other parts that are less likely to be reused are to be thrown out responsibly or recycled where possible. Special attention should be paid to proper disposal of hazardous materials.
